%author: Pedro Pinto; v4; 22-nov 2021
\documentclass[11pt,a4paper]{report}
\usepackage[utf8]{inputenc}
\usepackage[portuguese,english]{babel}
\usepackage{lipsum}
\usepackage{amsmath}
\usepackage{graphicx}
\usepackage{tabularx}
\usepackage{adjustbox}
%\usepackage[colorinlistoftodos]{todonotes}
\usepackage{csquotes}
\usepackage{comment}
\usepackage{imakeidx}
%tabela
\usepackage{multirow}
\usepackage{xcolor}
\usepackage[margin=3cm]{geometry}
\usepackage[hidelinks]{hyperref}
\usepackage[toc,acronym,nopostdot,nonumberlist]{glossaries}
\usepackage{titling}
\usepackage{tikzpagenodes}
\usepackage[ddmmyyyy]{datetime}
\usepackage{setspace}
\usepackage{indentfirst}
\usepackage[toc,page]{appendix}
%para definir a localização das tabelas e imagens em modo strict
%\usepackage{placeins}
\usepackage{biblatex}
\addbibresource{ref.bib}
\makeglossaries
\input{vars}
\doublespacing
%\pagestyle{headings}
\usepackage{fancyhdr}
\usepackage{lastpage}
\usepackage{pdfpages}
\usepackage{listings}
\definecolor{graygray}{gray}{0.95}
%\usepackage{helvet}
%\renewcommand{\familydefault}{\sfdefault}
\usepackage{graphicx}
\usepackage{eso-pic}
\usepackage{ragged2e}
\usepackage{changepage}
\usepackage[absolute,overlay]{textpos}

\newcommand\BackgroundPic{
    \put(0,0){
        \parbox[b][\paperheight]{\paperwidth}{
            \vfill
            \centering
            \includegraphics[width=\paperwidth,height=\paperheight,keepaspectratio]{Capa.pdf}%
            \vfill
        }
    }
}

\begin{document}
\thispagestyle{empty}
\ClearShipoutPicture
\begin{adjustwidth}{0pt}{20pt}
\begin{flushright}
\includegraphics[scale=0.45]{figs/logo-ipvc-old.pdf}\\
\vspace*{2cm}
\textsf{\huge Title of the Thesis/Project/Dissertation}\\
\vspace{3cm}
\large \textsf{Autor(a)}\\
\textcolor{gray}{\textsf{\Large Author}}\\
\vspace{3cm}
\vspace{1cm}
\large \textsf{Trabalho orientado por}\\
\textcolor{gray}{\large \textsf{Professor Name - E.g. Prof. John Doe}}\\
\vspace{2cm}
\textcolor{gray}{\large \textsf{Mestrado em Cibersegurança}}\\
\vspace{2cm}
\large \textsf{DD de MM de 20YY}\\
\end{flushright}
\end{adjustwidth}

\newpage
\ClearShipoutPicture
\begin{titlepage}
\begin{tikzpicture}[remember picture,overlay,shift={(current page.center)}]
\node[anchor=center,xshift=-2.8cm,yshift=10.4cm]{\includegraphics[scale=0.18]{figs/MCiber-logo.pdf}};
\node[anchor=center,xshift=2.8cm,yshift=10.3cm]{\includegraphics[scale=0.7]{figs/letter-MCiber.pdf}};
\node[anchor=center,xshift=7cm,yshift=-13cm]{\includegraphics[scale=0.6]{figs/MCiber-end2.pdf}};
\end{tikzpicture}

\centering
\vspace{4.5cm}
\huge Title of the\\Thesis\\
\vspace{1.5cm}
\large a master's thesis authored by\\
\Large Full Name of the student\\
\vspace{1.5cm}
\large and supervised by\\
\large FULL NAME\\\small Professor Assistente/Adjunto/Coordenador, IPVC\\
\vspace{0.2cm}
\large FULL NAME\\\small Professor Assistente/Adjunto/Coordenador, IPVC\\
\vspace{1.5cm}
\normalsize This thesis was submitted in partial fulfilment of the requirements for the\\Master's degree in Cybersecurity at the Instituto Politécnico de Viana do Castelo\\
\vspace{0.5cm}
\includegraphics[scale=0.65]{figs/logo-IPVC-s.pdf}\\
\vspace{0.2cm}
\newdateformat{daymonthyear}{\THEDAY\ of \monthname[\THEMONTH], \THEYEAR}
\daymonthyear\today \\
%December, 2019
%\large v2.3
\end{titlepage}


\input{acro_list}


%%%%%%%%%%%%%%%%%%%%%%%%%%%%%%%%%%%%%%%%%%%%%%%%%%%%%  

\selectlanguage{english}
\pagenumbering{roman}
\chapter*{\centering\Large\bfseries Abstract}
%context

%problem

%proposal

%solution

    \lipsum[1]
    
    \vspace{2cm}
    \textbf{Keywords:} word1. word2. word3. word4. 


\selectlanguage{portuguese}
\chapter*{\centering\Large\bfseries Resumo}
    \lipsum[1]
    
    \vspace{2cm}
    \textbf{Palavras-chave:} word1. word2. word3. word4. 

\selectlanguage{english}

\chapter*{Aknowledgements}

I want to thanks.. 


%insert index
\tableofcontents
\clearpage 
	\phantomsection
	\addcontentsline{toc}{chapter}{\listfigurename}			% The List of Figures (Do not modify)
	\listoffigures
\clearpage 	
	\phantomsection
	\addcontentsline{toc}{chapter}{\listtablename}			% The List of Tables (Do not modify)
	\listoftables
	
%in case there are listings with code:	
\clearpage
    \renewcommand\lstlistlistingname{List of Listings}
	\phantomsection
	\addcontentsline{toc}{chapter}{List of Listings}			% The List of Tables (Do not modify)
	\lstlistoflistings	



\glsaddall
%\printglossary[type=\acronymtype]
\printglossary[type=\acronymtype,title={List of Abbreviations},toctitle={List of Abbreviations}]
%%%%%%%%%%%%%%%%%%%%%%%%%%%%%%%%%%%%%%%%%%%%%%%%%%%%% 
\clearpage

\pagestyle{fancy}
\fancyhf{}
%\fancyhead[LE,LO]{Title Short}
\fancyhead[CO]{\nouppercase{\leftmark}}
%\fancyfoot[CE,CO]{\leftmark}
\fancyfoot[CO]{\thepage}
\cfoot{Page \thepage\ of \pageref{lastpage1}}
\renewcommand{\headrulewidth}{0.4pt}
\renewcommand{\footrulewidth}{0.4pt}
\fancypagestyle{plain}{%
  \fancyhf{}%
  \fancyfoot[C]{Page \thepage\ of \pageref{lastpage1}}%
  \renewcommand{\headrulewidth}{0pt}% Line at the header invisible
  \renewcommand{\footrulewidth}{0.4pt}% Line at the footer visible
}

\glsresetall
\chapter{Introduction}
\pagenumbering{arabic}
\label{chap:introduction}

In this Chapter ... 

IMPORTANT:
\begin{itemize}
    \item PLEASE check the file vars.tex: I'm using it in 10\kbps.
    \item Using glossaries for the first time: \gls{wn}
    \item Using glossaries after first time: \gls{wn}
    \item Using glossaries in plural: \glspl{wn}
    \item Using glossaries extended: \glsdesc{wn}
\end{itemize}

How to reference and insert the Fig.~\ref{fig:labeloffigure}. 

\begin{figure}[ht]
    \centering
    \includegraphics[width=0.3\textwidth]{figs/MCiber-logo.pdf}
    \caption{This is a Figure}
    \label{fig:labeloffigure}
\end{figure}

This is an example of a citation \cite{Pedro2019}.   

This is a way for you to generate tables:

\url{https://www.tablesgenerator.com/}

Table~\ref{table-x} presents an example.

\begin{table}[h!]
\caption{This is a Table}
\center
\begin{tabular}{|l|c|c|}
\hline
 & Test A & Test B \\ \hline
P1 & 13,5 & 5,1 \\ \hline
P2 & 12,4 & 6,2 \\ \hline
P3 & 10,2 & 7,2 \\ \hline
P4 & 9,5 & 10,1 \\ \hline
\end{tabular}
\label{table-x}
\end{table}

In Listing~\ref{lst1} is presented a code. Other example in Listing~\ref{lst2}.

\begin{lstlisting}[xleftmargin = 1cm,basicstyle=\footnotesize,
    numbers=left,captionpos=b,caption=Script for a basic Hello World, label=lst1, language=Python, backgroundcolor = \color{graygray}, breaklines]
{
    # This program prints Hello, world!

    print('Hello, world!')
}
\end{lstlisting}


\begin{lstlisting}[xleftmargin = 1cm, basicstyle=\small\tt,captionpos=b,caption=Script for a basic Hello World, label=lst2, frame=single, numbers=left, language=Python]
{
    # This program prints Hello, world!

    print('Hello, world!')
}
\end{lstlisting}


\lipsum[1]


\section{Context} 
\label{sec:context}
%Context
\lipsum[1] 

%context focus
\lipsum[1] 

\section{Problem Statement and Motivation}
\label{sec:motivation}
%problem to solve
\lipsum[1] 

%point the motivation
\lipsum[1]

\section{Objectives}
\label{sec:obj}
\lipsum[1]

\section{Contributions}
\label{sec:cont}
The contributions provides the following contributions:
\begin{enumerate}
    \item \lipsum[1]
    \item \lipsum[1]
\end{enumerate}

\section{Organization}
\label{sec:org}
\lipsum[1] 

\chapter{Title of the Chapter}
\label{cap:name1}

\lipsum[1] 

\section{Section 1}
\label{sec:sec1}

\lipsum[1] 

\begin{itemize}
    \item \lipsum[1]
    \item \lipsum[1]
\end{itemize}

%%%%%%%%%%%%%%%%%%%%%%%%%%%%%%%%%%%%%%%%%%%%%%%%%%%%%  
\chapter{Title of the Chapter}
\label{cap:name2}

\lipsum[1] 

\section{Section 2}
\label{sec:sec2}

\lipsum[1]

\chapter{Conclusions}
\label{cap:conc}
\lipsum[1]

\clearpage
\addcontentsline{toc}{chapter}{References}
\printbibliography[title={References}]

\label{lastpage1}
\clearpage

\pagestyle{fancy}
\fancyhf{}
%\fancyhead[LE,LO]{Title Short}
\fancyhead[RO]{\nouppercase{\leftmark}}
%\fancyfoot[CE,CO]{\leftmark}
\fancyfoot[CO]{\thepage}
\cfoot{Page \thepage\ of \pageref{LastPage2}}
\renewcommand{\headrulewidth}{0.4pt}
\renewcommand{\footrulewidth}{0.4pt}
\fancypagestyle{plain}{%
  \fancyhf{}%
  \fancyfoot[C]{Page \thepage\ of \pageref{lastpage2}}%
  \renewcommand{\headrulewidth}{0pt}% Line at the header invisible
  \renewcommand{\footrulewidth}{0.4pt}% Line at the footer visible
}
\pagenumbering{arabic}% resets `page` counter to 1
\renewcommand*{\thepage}{A\arabic{page}}
\appendix
\begin{appendices}
\chapter{Title of Appendix}
The contents...


\chapter{Title of other Appendix}
The contents...
\end{appendices}
\label{lastpage2}
\end{document}


%HELP
%https://tikz.dev/pgfplots/reference-2dplots